\section{One character and compatible measures on the actual setting}
\label{sec:global-measure-rescaling-v53}

The purpose of this chapter is the original \texttt{RescaleStep3}.
We start with the supplied Hida setting, its supplied adelic bridge, and
its already normalized measures at infinity.  We construct the fiber
measures which determine the global stabilizer measures and the local
flag-stabilizer measures.  We then use AMD13--AMD17 to fill in the finite
factors without changing those measures at infinity.  The character is
chosen once, before any word, finite place, or auxiliary tuple.

The construction of a fiber measure is different from the assertion
that the volume of an arithmetic quotient equals a Siegel--Weil
constant.  Here the normalization is fixed by local Fourier inversion
and local disintegration.  At almost every finite place we compute its
mass exactly.  The resulting positive, convergent product constructs the
global fiber measure and supplies its nonvanishing.  Consequently the
measure construction below does not require a theta-integral identity.
We state the original source endpoint with its exact
\texttt{SiegelWeilUnitaryConvergent} premise; that larger premise, and the
later arithmetic volume comparisons which use it, are not proved in this
chapter.

\subsection{Notation and the precise measure contracts}

Let $E/F$ be a CM extension, with involution $c$, and let
$h(x,y)=x^*Jy$ be a nondegenerate Hermitian form on $E^n$.
For a place $v$ put $k_v=F_v$, $A_v=E\otimes_F k_v$ and
$G_v=U(J)(k_v)$.  The algebra $A_v$ is either a quadratic field or
$k_v\times k_v$; at infinity it is $\mathbb C$ over $\mathbb R$.
All assertions below include the split algebra.  We use column notation
for an ordered tuple, so
\[
 \begin{gathered}
 V_{v,a}=\operatorname{Mat}_{n\times a}(A_v),\qquad
 \gamma_v(X)=X^*JX,\\
 U_{v,a}=\{X:\text{the columns of }X\text{ are }A_v\text{-independent}\}.
 \end{gathered}
 \tag{GMR.1}
\]
For the product of the infinite places the notation has the same
meaning, component by component.  On a split algebra, independence
means independence in both factors.  Write $\mathcal H_{v,a}$ for the
$k_v$-space of Hermitian $a\times a$ matrices.  Its dimension is $a^2$;
$V_{v,a}$ has dimension $2na$ over $k_v$.

An adelic character $\psi$ satisfying \texttt{IsCharAF} is continuous,
unitary, trivial on the diagonal copy of $F$, and nontrivial.  The
normalization \texttt{IsStdChar} is exactly
\[
 \psi_\infty(x)=\exp\left(2\pi i\sum_{\tau:F\hookrightarrow\mathbb R}x_\tau\right).
 \tag{GMR.2}
\]
The pairings defining the source measures \texttt{SdW} and
\texttt{SdHerm} are the trace pairings \texttt{BR} and
\texttt{pairHerm}, with this character.  We always use those pairings;
an arbitrary normalization of coordinate Lebesgue or additive Haar
measure is not substituted for a self-dual measure.

Here is the full fiber-family condition used below.  Given self-dual
measures $dX$ on $V_{v,a}$ and $dT$ on $\mathcal H_{v,a}$, a family
$\omega_{v,T}$ satisfies \texttt{IsFibreFamily} when
\[
\begin{split}
 &\omega_{v,T}\bigl(V_{v,a}\setminus(\gamma_v^{-1}(T)\cap U_{v,a})\bigr)=0,\\
 &\int_{U_{v,a}} f(X)\,dX
       =\int_{\mathcal H_{v,a}}\int f(X)\,d\omega_{v,T}(X)\,dT
       \quad(f\geq0\text{ measurable}),
\end{split}
 \tag{GMR.3}
\]
and, for every real continuous compactly supported $f$ whose support is
contained in $U_{v,a}$, the fiber integral is finite and continuous in
\emph{every} $T\in\mathcal H_{v,a}$.  There is no exclusion of singular
Hermitian matrices $T$ in this condition.  Measures are understood on
the ambient tuple space, by extension by zero from the independent
locus.

For a rational nondegenerate Hermitian matrix $T$, the original
\texttt{IsOmT} requires a measure $\nu_T$ supported on the adelic level
set $\gamma=T$ and the following assertion for \emph{all} self-dual
measures and \emph{all} fiber families satisfying (GMR.3).
If
\[
 \Phi(X)=\Phi_\infty(X_\infty)\Phi_f(X_f),\qquad
 \Phi_f(X_f)=\prod_v\Phi_v(X_v),
 \tag{GMR.4}
\]
where $\Phi_\infty$ is continuous with compact support, every $\Phi_v$
is locally constant with compact support, and $\Phi_v$ is the indicator
of the coordinate lattice attached to a single rational $F$-basis of
$E$ at all but finitely many places, then $\Phi$ is integrable,
the product below is multipliable, and
\[
 \int\Phi\,d\nu_T
   =\left(\int\Phi_\infty\,d\omega_{\infty,T}\right)
      \prod_v\int\Phi_v\,d\omega_{v,T}.
 \tag{GMR.5}
\]
This includes the existence quantifiers in the source's
\texttt{AllSome}: the self-dual measures and the fiber families must
exist, and the conclusion must hold for every permitted choice.

For an actual tuple $X_0$ and its actual stabilizer $H$, the condition
\texttt{IsCompatT}, or \texttt{IsCompatLoc} at one place, says that
$\nu_T$ is supported on the actual orbit and
\[
 \int_G f(g)\,d\mu_G(g)
  =\int_{G X_0}\int_H f(s(y)h)\,d\mu_H(h)\,d\nu_T(y)
 \tag{GMR.6}
\]
for every measurable $f:G\to[0,\infty]$ and \emph{every} section $s$
with $s(y)X_0=y$ on that orbit.  The original statement does not assume
that this section is measurable.  We prove precisely this version by
showing that the inner integral is independent of the section and has
a measurable representative defined by local sections.

The principal source interfaces are
\texttt{VolFromSW.IsOmT}, \texttt{VolFromSW.IsCompatT},
\texttt{LocVol.IsOm}, \texttt{LocVol.LocTieW}, and
\texttt{SettingExistsAdapted.RescaleStep3}.  Their literal definitions
and the producer bodies are included in the source manifest.  AMD1,
AMD4--AMD8 and AMD13--AMD17 are used at their reviewed original scopes.

% BEGIN GMR1
\paragraph{GMR1. A single standard character and its local conductors.}
\label{claim:GMR1}
For every totally real number field $F$ there is a unique character
$\psi^0$ satisfying \texttt{IsCharAF} and (GMR.2).
Every character satisfying \texttt{IsCharAF} has the form
$\psi(x)=\psi^0(bx)$ for a unique $b\in F^\times$.
In particular its restriction to each completion is nontrivial, and
its finite-place conductor is $\mathcal O_v$ at all but finitely many
places.  These conclusions are made before any tuple or word is chosen.

\begin{proof}
On $\mathbb A_{\mathbb Q}$ take the character whose real factor is
$e^{2\pi ix}$ and whose $p$-adic factor is
$e^{-2\pi i\{x\}_p}$, where $\{x\}_p$ is the finite negative-power
part of the $p$-adic expansion, viewed in $\mathbb Q/\mathbb Z$.
The sum of the negative-power parts of a rational number differs from
that number by an integer.  Thus the product of these factors is one
on $\mathbb Q$.  Compose with the adelic trace from $F$ to $\mathbb Q$.
The local traces over the places above $p$ add to the rational trace;
the infinite trace is the sum over the real embeddings.  We obtain
(GMR.2), triviality on $F$, continuity, and absolute value one.
The infinite factor is nontrivial, so the whole character is nontrivial.
The annihilator of $\mathcal O_v$ for its local trace pairing is the
inverse different.  Outside the finitely many ramified rational primes
this is $\mathcal O_v$ itself.

We recall the character argument which also verifies that no local
factor can disappear for another permitted $\psi$.  A continuous
character of a real vector space is the exponential of a real linear
functional.  A continuous character of a nonarchimedean local field is
of the form $x\mapsto\psi^0_v(b_vx)$: on a compact open additive subgroup
it factors through a finite quotient, whose character group is its
trace dual, and the compatible characters on the increasing groups
$\varpi^{-m}\mathcal O_v$ give the unique parameter $b_v$.
Continuity on the finite adeles implies triviality on
$\prod_{v\notin S}\mathcal O_v$ for some finite $S$.  One way to see
this is to choose a small arc about one containing no nontrivial
subgroup of the unit circle and use continuity on that compact product.
Consequently the parameters $b_v$ form an adele $b$.

The annihilator of diagonal $F$ for $\psi^0$ is diagonal $F$ itself.
For completeness, choose $r\in F$ with $b_f-r\in\prod_v\mathcal O_v$;
this is the Chinese remainder theorem applied to the finitely many
denominators of $b_f$.  Triviality on $\mathcal O_F$ now says that
$b_\infty-r$ lies in the trace-dual lattice of $\mathcal O_F$ in
$F_\infty$.  That lattice is the image of the inverse different, so
$b_\infty-r$ is the infinite image of an element $d\in F$.
Subtract $r+d$.  The remaining parameter has zero infinite component,
and its finite character is trivial on $F$.  The Chinese remainder
theorem also proves that $F$ is dense in $\mathbb A_{F,f}$, so this finite
character is identically one.  Local nondegeneracy of the trace pairing
then makes every remaining finite parameter zero.  This proves the
annihilator assertion and hence $b\in F$.  Nontriviality of $\psi$ gives
$b\ne0$.  At almost every place $b$ is a unit, so the conductor assertion
follows from the assertion for $\psi^0$.

Finally, if the infinite factor is (GMR.2), the infinite images of $b$
are all one, so $b=1$.  This proves uniqueness of the standard
character, as well as the original \texttt{standard\_character\_exists}
assertion.
\end{proof}
% END GMR1

% BEGIN GMR2
\paragraph{GMR2. Self-dual measures exist and are uniquely normalized.}
\label{claim:GMR2}
For every nontrivial local factor supplied by GMR1, the source pairings
on $V_{v,a}$, on $\mathcal H_{v,a}$, and on $A_v^I$ with pairing
$\psi_v(\sum_{i\in I}\operatorname{Tr}_{A_v/k_v}(x_i y_i))$
have unique self-dual additive Haar measures.  The same holds for the
product of the infinite places.  These measures satisfy the literal
source Fourier-inversion predicate, not just invariance under addition.
Whenever an integral lattice is its own annihilator, its self-dual mass
is one.

\begin{proof}
Each pairing is nondegenerate over the base field.  For tuple space
this follows from invertibility of $J$ and nondegeneracy of the trace
pairing of the quadratic etale algebra.  For Hermitian matrices,
diagonal entries give the scalar trace pairing and each pair of
off-diagonal entries gives the trace pairing of $A_v$; the coefficients
which occur are nonzero in characteristic zero.  This also proves the
assertion for a split algebra.  The direct sum pairing on $A_v^I$ is
nondegenerate, including $I=\varnothing$.

Here is the normalization argument.  In real coordinates, Fourier
inversion for $e^{2\pi i x\cdot y}$ uses ordinary Lebesgue measure.
Replacing this pairing by $e^{2\pi i x^tBy}$ multiplies the self-dual
measure by $|\det B|^{1/2}$.  Indeed a change of variable in the two
successive transforms contributes $|\det B|^{-1}$, exactly cancelling
the square of that factor.  The same calculation applies over a local
field.  It can also be obtained directly from a compact open lattice
$L$: the transform of $1_L$ is
\[
 \widehat{1_L}(y)=\mu(L)1_{L^\perp}(y),\qquad
 \widehat{\widehat{1_L}}(z)=\mu(L)\mu(L^\perp)1_{-L}(z).
 \tag{GMR.7}
\]
Starting with any additive Haar measure, there is exactly one positive
rescaling for which $\mu(L)\mu(L^\perp)=1$.  Fourier inversion follows
first on locally constant compactly supported functions by finite
Fourier inversion on lattice quotients, then on continuous compactly
supported functions whenever the transform is integrable.  In the real
case the corresponding extension follows by convolution with Gaussian
approximate identities and dominated convergence for the integrable
transform.  This is precisely the conditional inversion assertion in
the source definition of \texttt{SelfDual}.

To check uniqueness even for that conditional assertion, two additive
Haar measures differ by a positive factor $c$.  Apply the asserted
inversion to a nonzero compactly supported locally constant function
at a finite place, or to a smooth compactly supported function at
infinity.  Its Fourier transform is integrable; two successive
transforms for $c\mu$ multiply the answer for $\mu$ by $c^2$.
Thus $c^2=1$, and positivity gives $c=1$.  Formula (GMR.7) now gives
$\mu(L)=1$ when $L=L^\perp$.  Products of the individual infinite-place
measures give the asserted measure at infinity by Fubini and the same
uniqueness argument.
\end{proof}
% END GMR2

% BEGIN GMR3
\paragraph{GMR3. The Gram map is a submersion on independent tuples.}
\label{claim:GMR3}
Over every $k_v$ and at the product of the infinite places, the map
$\gamma:U_{v,a}\to\mathcal H_{v,a}$ is an analytic submersion.
This remains true when $\gamma(X)$ is singular.  Near every independent
$X$, the fiber has dimension $2na-a^2$ over the base field.

\begin{proof}
Since the columns of $X$ are independent and $J$ is invertible, the
map $X^*J:A_v^n\to A_v^a$ is onto.  Choose a right inverse $W$, so
$X^*JW=1$.  For a Hermitian matrix $S$, put $Y=WS/2$.
The differential of the Gram map is
\[
 D\gamma_X(Y)=X^*JY+Y^*JX=S/2+S^*/2=S.
 \tag{GMR.8}
\]
All divisions are in a characteristic-zero field; the argument also
works at places above two.  Over a split algebra the right inverse
is chosen separately in each factor.  Over a product of infinite
completions it is chosen componentwise.  Thus the differential is
surjective as a map of base-field vector spaces.

Choose $a^2$ of the real or local-field coordinates on which this
differential is invertible and retain the remaining coordinates $z$.
The inverse function theorem applied to $X\mapsto(z(X),\gamma(X))$
gives local analytic coordinates $(z,T)$ in which $\gamma$ is projection
onto $T$.  There are $2na-a^2$ coordinates $z$.
Nothing in this argument used $\det\gamma(X)\ne0$; independence is the
necessary hypothesis.
\end{proof}
% END GMR3

% BEGIN GMR4
\paragraph{GMR4. Construction of the full continuous fiber family.}
\label{claim:GMR4}
With the self-dual measures of GMR2 there is a family
$\omega_{v,T}$ satisfying all of (GMR.3), with continuity on the whole
Hermitian space.  On the independent part of every fiber it is a Radon
measure, positive on every nonempty relatively open set.  It is zero
when that independent fiber is empty.  The same assertion holds at
infinity.  More generally, the assertion holds for any prescribed
pair of additive Haar measures on tuple space and Hermitian space;
self-duality is not needed for this local construction.
In particular the source's chosen \texttt{dX}, \texttt{dT}
and \texttt{omFam} satisfy their defining alternatives rather than
their zero fallback alternatives.

\begin{proof}
Use a chart from GMR3 and write its inverse as $X=\chi(z,T)$.
Choose coordinate Haar measures $dz$ and $dT$ and express the ambient
measure in this chart as
\[
 dX=j(z,T)\,dz\,dT,\qquad j(z,T)>0.
 \tag{GMR.9}
\]
Here $j$ is the absolute Jacobian, including the positive constants
relating the coordinate measures to the prescribed additive Haar
measures.  It is continuous; over a nonarchimedean field the absolute
value of a nonvanishing analytic determinant is locally constant.
On the part of the fiber in this chart define its measure by
$j(z,T)\,dz$.  On overlapping charts the coordinate change fixes $T$.
Its derivative is block triangular, so its full Jacobian is its
vertical Jacobian.  The change-of-variables formula therefore says
that the two fiber measures agree on the overlap, for each $T$.

Take a countable locally finite chart cover of $U_{v,a}$ and glue these
agreeing measures.  This is an ordinary gluing of measures on open
sets: one can disjointify a countable cover and sum the restrictions,
and agreement on overlaps makes the result independent of the
disjointification.  Extend each resulting measure by zero to the
ambient tuple space.  The support assertion in (GMR.3) is immediate.
Fubini in each chart, followed by a countable disjointification and
monotone convergence, proves the disintegration identity for every
nonnegative measurable function.  This construction also shows that
the fiber integrals occurring there are measurable in $T$.

Let $f$ be continuous with compact support contained in $U_{v,a}$.
Cover its support by finitely many such charts and take a continuous
partition of unity subordinate to slightly smaller charts.  In a
nonarchimedean space a finite clopen refinement gives the same
partition.  Each resulting fiber integral is an integral of a
continuous function with a common compact vertical support, locally
in $T$.  Its absolute value is bounded by a fixed integrable function
because $j$ is bounded on that compact set.  Dominated convergence
proves continuity and finiteness.  The summands extend by zero past
the boundaries of their smaller charts, so this proves continuity on
all of $\mathcal H_{v,a}$, including parameters with empty fibers.
The construction is Radon on the fiber inside $U_{v,a}$ and its chart
density is strictly positive.  These observations prove the remaining
assertions.  If a singular fiber approaches a dependent tuple, no
unproved assertion of ambient local finiteness is being made: the
Radon assertion here is on the independent fiber, exactly where the
quotient argument below takes place.
\end{proof}
% END GMR4

% BEGIN GMR5
\paragraph{GMR5. Pointwise uniqueness and covariance of the fibers.}
\label{claim:GMR5}
For any fixed pair of additive Haar measures, the family in GMR4 is
the unique family satisfying the original \texttt{IsFibreFamily} condition.
In particular uniqueness holds at each singular Hermitian parameter,
not merely almost everywhere.  Every unitary left action preserves
$\omega_{v,T}$.  If $u$ is upper unitriangular, right multiplication
$X\mapsto Xu$ carries $\omega_{v,T}$ to $\omega_{v,u^*Tu}$.
The same assertions hold at infinity.

\begin{proof}
For a real continuous $f$ of compact support in the independent locus,
two proposed families have fiber integrals with equal integrals over
every measurable set of Hermitian parameters, by disintegration.
The functions are consequently equal almost everywhere.  Both are
continuous by the defining condition, and additive Haar measure is
positive on every nonempty open set.  They are therefore equal
everywhere.  For each fixed $T$, these test functions determine the
measure on the independent fiber: compact subsets have continuous
cutoffs with compact support in the independent locus, and the
defining integrability gives local finiteness there.  The measures are
Radon on that locally compact space, and a compact exhaustion followed
by monotone convergence determines their values on all Borel sets.
Both vanish outside that locus.  Hence the entire ambient measures
are equal at every $T$.

More generally, suppose invertible linear maps $A$ on tuple space and
$C$ on Hermitian space satisfy $\gamma A=C\gamma$ and preserve the
independent locus.  Pushing (GMR.3) through these changes of variables
and using uniqueness gives
\[
 A_*\omega_T=\frac{|\det C|}{|\det A|}\,\omega_{CT}.
 \tag{GMR.10}
\]
The determinant is over the base local field, with the product of
absolute values at infinity.  For $A(X)=gX$, $g\in U(J)$, its determinant
is $\operatorname{Nm}_{A_v/k_v}(\det g)^a=1$ and $C=1$.
For $A(X)=Xu$ with $u$ upper unitriangular, both $A$ and
$C(T)=u^*Tu$ are unipotent linear maps and have determinant one.
Thus (GMR.10) gives both claimed covariance assertions.  The existence
and uniqueness in GMR2 also show that these conclusions apply to
\emph{any} source witnesses for the self-dual measures and fiber
families, including the witnesses chosen by \texttt{omFam}.
\end{proof}
% END GMR5

% BEGIN GMR6
\paragraph{GMR6. Exact counts of nondegenerate Hermitian frames.}
\label{claim:GMR6}
Let $k=\mathbb F_q$ have odd characteristic, let $A$ be its quadratic
etale algebra with its nontrivial involution, and let $J,T$ be
nondegenerate Hermitian forms of ranks $n,a$, where $0\leq a\leq n$.
Put $d=2na-a^2$, and put $\varepsilon=1$ in the split case and
$\varepsilon=-1$ in the field case.  Then
\[
 \#\{X\in\operatorname{Mat}_{n\times a}(A):X^*JX=T\}
  =q^d\prod_{j=n-a+1}^{n}(1-\varepsilon^j q^{-j}).
 \tag{GMR.11}
\]
This gives, in particular, the full original finite-field count bound
\texttt{FFCount} when $a+2\leq n$.

\begin{proof}
Over $\mathbb F_{q^2}$ every nondegenerate Hermitian form is equivalent
to the identity: diagonalize by orthogonal elimination and rescale
each nonzero diagonal entry, using the surjectivity of the norm to
$\mathbb F_q^\times$.  The unitary group acts transitively on the
ordered frames with Gram matrix $T$, since an isometry on their
nondegenerate spans extends by choosing orthonormal bases of their
orthogonal complements.  The stabilizer fixes the span pointwise and
is $U_{n-a}(q)$.

Here is the order calculation.  For a nontrivial additive character
$\eta$ of $\mathbb F_q$ and $s\ne0$,
\[
 \sum_{z\in\mathbb F_{q^2}}\eta(s z\bar z)
   =1+(q+1)\sum_{t\in\mathbb F_q^\times}\eta(st)=-q.
\]
Character orthogonality therefore gives
\[
 \#\{z\in(\mathbb F_{q^2})^m:\textstyle\sum z_i\bar z_i=1\}
  =q^{2m-1}-(-1)^m q^{m-1}.
\]
Choose the first unit vector of an orthonormal basis and then an
orthonormal basis of its complement.  Induction gives
\[
 |U_m(q)|=q^{m^2}\prod_{j=1}^{m}(1-(-1)^j q^{-j}).
 \tag{GMR.12}
\]
Taking $|U_n|/|U_{n-a}|$ proves (GMR.11) in the field case.

For $A=k\times k$, a Hermitian space is a perfect pairing between two
$k$-spaces.  Its unitary group is $\operatorname{GL}_n(k)$, and the
stabilizer of a nondegenerate ordered frame is
$\operatorname{GL}_{n-a}(k)$.  Counting successively independent
columns gives $|\operatorname{GL}_m(k)|=q^{m^2}\prod_{j=1}^m(1-q^{-j})$.
Its quotient gives the split formula.  Equivalently one may choose
the first-component frame in $\prod_{i=0}^{a-1}(q^n-q^i)$ ways and
then choose its paired second-component frame in $q^{a(n-a)}$ ways;
this is exactly the source's \texttt{FFSplit} formula.

To obtain the uniform bound, set $s=n-a+1$.  When $a+2\leq n$, $s\geq3$.
For $b_j=-\varepsilon^j q^{-j}$,
$\sum_{j=s}^n|b_j|\leq2q^{-s}$ and
$|\prod(1+b_j)-1|\leq\exp(\sum|b_j|)-1\leq4q^{-s}$.
Thus (GMR.11) differs from $q^d$ by at most $4q^{d-2}$.
The constant $4$, including the empty product $a=0$, supplies the
existential constant in the full original \texttt{FFCount}.
\end{proof}
% END GMR6

% BEGIN GMR7
\paragraph{GMR7. Exact good-place mass of the canonical Gram fiber.}
\label{claim:GMR7}
Fix $1\leq a$, $2a<n$, a rational tuple $X_0$ with nondegenerate
$T=X_0^*JX_0$, an adelic character satisfying \texttt{IsCharAF}, and
any rational $F$-basis used to define the source lattices.
Outside a finite set of places, the compact fiber
\[
 C_v=\{X\in\operatorname{Mat}_{n\times a}(\mathcal O_{E,v}):X^*JX=T\}
\]
is nonempty and has canonical fiber mass
\[
 a_v:=\omega_{v,T}(C_v)
       =\prod_{j=n-a+1}^{n}(1-\varepsilon_v^j q_v^{-j}),
 \tag{GMR.13}
\]
where $\varepsilon_v=1$ at split places and $-1$ at inert unramified
places.  Formula (GMR.13) holds for every permitted self-dual measure
and fiber family.

\begin{proof}
Exclude places above two, places ramified in $E/F$ or in the different
of $F$, nonunit places of the character parameter in GMR1, denominators
of $J,T,X_0$ and their nonzero determinants, and the finitely many
places where the chosen rational basis lattice is not
$\mathcal O_{E,v}$.  At a remaining place $A_v$ is unramified etale,
$J,T$ are integral and unimodular, and the character has conductor
$\mathcal O_v$.  The two trace-pairing lattices on tuple space and
Hermitian space are perfect integral duals.  Thus their self-dual
masses are one by GMR2.  The specified tuple $X_0$ belongs to $C_v$.

Reduction modulo $\varpi$ of any solution has independent columns,
because its Gram matrix $\bar T$ is invertible.  At such a solution
the differential (GMR.8) is onto over the residue field.  An arbitrary
lift of a solution modulo $\varpi^m$ has a Gram error at order
$\varpi^m$; correcting it by $\varpi^mY$ changes that error by
$D\gamma(Y)$ modulo $\varpi^{m+1}$.  Surjectivity shows that each
solution has exactly $q_v^{2na-a^2}$ lifts to the next level.
If $N_m$ is the number of solutions modulo $\varpi^m$, then
$N_m=N_1q_v^{(2na-a^2)(m-1)}$.

The ambient volume of the integral tuples satisfying
$\gamma(X)\equiv T\pmod{\varpi^m}$ is $N_mq_v^{-2nam}$.
The volume of the corresponding Hermitian ball is $q_v^{-a^2m}$.
Their quotient is consequently $N_1q_v^{-(2na-a^2)}$, independently
of $m\geq1$.  By (GMR.3) this quotient is the average of the fiber
mass of the integral tuples over that ball.  To apply the continuity
clause literally, multiply the lattice indicator by the indicator
of $\gamma(X)\equiv T\pmod\varpi$: its support is a compact subset
of the independent locus, and on all these balls it has the same
fiber integral.  Continuity from GMR4 lets $m$ tend to infinity and
identifies the average with $\omega_{v,T}(C_v)$.
Now GMR6 gives exactly (GMR.13).  GMR2 and GMR5 remove all dependence
on the chosen witnesses.  Enlarging the exceptional set handles any
other rational basis, which is precisely the basis quantifier in
\texttt{IsPureFin}.
\end{proof}
% END GMR7

% BEGIN GMR8
\paragraph{GMR8. The density product is convergent and strictly positive.}
\label{claim:GMR8}
In the range of GMR7, put $s=n-a+1$.  For every good place,
\[
 a_v>0,\qquad |a_v-1|\leq4q_v^{-s},\qquad
 |\log a_v|\leq4q_v^{-s},\qquad s\geq a+2\geq3.
 \tag{GMR.14}
\]
Hence every product of the $a_v$ outside a fixed finite exceptional
set converges to a finite strictly positive real number.

\begin{proof}
Each factor in (GMR.13) is positive.  The bound for $a_v-1$ was proved
in GMR6.  Since $q_v^{-j}\leq1/8$ for $j\geq s\geq3$, the elementary
inequality $|\log(1+t)|\leq2|t|$ for $|t|\leq1/2$ gives
\[
 |\log a_v|\leq2\sum_{j=s}^nq_v^{-j}\leq4q_v^{-s}.
\]
There are at most $[F:\mathbb Q]$ places above a rational prime $p$,
and each has $q_v\geq p$.  Therefore
\[
 \sum_vq_v^{-s}\leq[F:\mathbb Q]\sum_{m=2}^{\infty}m^{-s}<\infty.
\]
The sum of the real logarithms converges absolutely; its exponential
is the desired finite positive product.  This proves positivity of
the normalization, rather than inferring it from convergence of a
possibly zero Euler product.
\end{proof}
% END GMR8

% BEGIN GMR9
\paragraph{GMR9. The global fiber measure satisfies the full \texttt{IsOmT}.}
\label{claim:GMR9}
Under the original range $1\leq a$, $2a<n$ and for the actual rational
nondegenerate Gram matrix $T$ of $X_0$, there is a nonzero Radon measure
$\nu_T$ on adelic tuple space satisfying the full original
\texttt{IsOmT}, including all basis, test-function, self-dual-measure
and fiber-family quantifiers in (GMR.5).  It is positive on every
nonempty relatively open subset of the adelic Gram fiber.

\begin{proof}
Write $Y_v=\gamma_v^{-1}(T)$ and $Y_\infty=\gamma_\infty^{-1}(T)$.
Nondegeneracy of $T$ implies that every tuple in these fibers is
independent.  Thus each $Y_v$ is closed in tuple space and its measure
$\omega_{v,T}$ is Radon by GMR4.  Choose an exceptional finite set
$S_0$ as in GMR7.  For $v\notin S_0$, the set $C_v\subset Y_v$ is
compact open of positive finite mass $a_v$.  Let
$\bar\omega_v=a_v^{-1}\omega_{v,T}$, so
$\bar\omega_v(C_v)=1$.

Here is the restricted-product measure construction with its scalar
displayed.  On
\[
 Y(S)=\prod_{v\in S}Y_v\times\prod_{v\notin S}C_v,
       \qquad S_0\subset S\text{ finite},
\]
use the product of $\omega_{v,T}$ for $v\in S_0$ and
$\bar\omega_v$ for $v\notin S_0$.  The infinitely many remaining
factors are probability measures on compact metrizable spaces, so
their product is the countably additive product measure.  When $S$
is enlarged, restriction to the smaller $Y(S)$ agrees, because the
new compact factors have normalized mass one.  These measures
therefore glue on the union $Y_f=\bigcup_SY(S)$.  Multiply the result
by the single positive number
\[
 A_0=\prod_{v\notin S_0}a_v.
 \tag{GMR.15}
\]
It exists by GMR8.  Call the resulting measure $\nu_{T,f}$.
Equivalently its restriction to $Y(S)$ is
\[
 \left(\prod_{v\notin S}a_v\right)
 \left(\bigotimes_{v\in S}\omega_{v,T}\right)
 \otimes\left(\bigotimes_{v\notin S}\bar\omega_v|_{C_v}\right).
 \tag{GMR.16}
\]
This formula proves directly that changing $S_0$ has no effect.

The usual restricted-product topology makes every $Y(S)$ open.
A compact set is covered by finitely many of these opens and hence
is contained in one of them.  Its projection at every place in $S$
is compact, so (GMR.16) proves finite mass on compact sets.  The
measure is Radon, either by regularity of these countable products
of compact metric probability spaces and local Radon factors, or by
compact exhaustion of each of their countably many charts.
It is nonzero and positive on every nonempty basic open set: the
finitely many specified local open sets have positive mass by GMR4,
and the remaining compact factors have mass one before multiplication
by a strictly positive scalar.

AMD1 identifies finite-adelic tensor coordinates with their
restricted-product local coordinates.  Under this identification
$Y_f$ is exactly the finite-adelic Gram fiber.  Set
$\nu_T=\omega_{\infty,T}\otimes\nu_{T,f}$ and push it into adelic
tuple space.  It is supported on $\gamma=T$, is Radon, and is nonzero
because the rational tuple $X_0$ supplies a point in every local
fiber and hence a nonempty product neighborhood.

Now take any test data in (GMR.4), using any rational basis.  Enlarge
$S$ to include $S_0$, the exceptional basis places, and the finitely
many places where $\Phi_v$ is not the standard lattice indicator.
On the fiber, $\Phi$ vanishes outside $Y_\infty\times Y(S)$, and its
support is compact in the infinite coordinate and the finitely many
unrestricted coordinates.  Formula (GMR.16) and Fubini therefore
prove absolute integrability and
\[
 \int\Phi\,d\nu_T
 =\left(\int\Phi_\infty\,d\omega_{\infty,T}\right)
   \left(\prod_{v\in S}\int\Phi_v\,d\omega_{v,T}\right)
   \prod_{v\notin S}a_v.
 \tag{GMR.17}
\]
The infinite product on the right is convergent by GMR8.  This also
proves multipliability when a finite local integral is zero or
complex; outside $S$ the factors are still the positive real $a_v$.
It is (GMR.5).  Finally, GMR2 and GMR5 prove equality with every
self-dual measure and every fiber family appearing in the source
\texttt{AllSome} quantifiers.  They also provide their existence.
Thus the full original contract, rather than just a factorization
for one selected set of witnesses, has been proved.
\end{proof}
% END GMR9

% BEGIN GMR10
\paragraph{GMR10. The adelic Gram fiber is the actual unitary orbit.}
\label{claim:GMR10}
For the data of GMR9, the full adelic Gram fiber equals
$U(J)(\mathbb A_F)X_0$.  The orbit map identifies the quotient by the
actual stabilizer with this fiber, has continuous local sections, and
is open.  The measure $\nu_T$ of GMR9 is invariant under the actual
adelic unitary group.

\begin{proof}
We give the needed extension argument rather than assume an orbit
identification as a measure input.  Over a quadratic field extension,
two nonisotropic vectors $u,v$ of the same norm $t$ can be matched by
a unitary transformation.  Set $w=u-v$, $d=h(w,u)$.  If $d\ne0$, the
map
\[
 R(x)=x-wd^{-1}h(w,x)
 \tag{GMR.18}
\]
sends $u$ to $v$ and is unitary, since
$h(w,w)=d+\bar d$ makes the two cross terms in
$h(Rx,Ry)-h(x,y)$ cancel its quadratic term.
If $d=0$, apply the same construction to $u,-v$, for which the
corresponding denominator is $2t\ne0$, and then negate on the current
orthogonal complement.  Working successively in the orthogonal
complements of already matched nonisotropic vectors extends every
isometry of nondegenerate subspaces.

The same argument gives extension on degenerate subspaces when it is
needed below.  Split a subspace into its radical and a nondegenerate
complement and first match the latter.  For a basis $e_i$ of the
remaining totally isotropic subspace, choose dual vectors $f_i$ with
$h(e_i,f_j)=\delta_{ij}$.  Replace the column $f$ by
$f-e\,h(f,f)/2$; its Gram matrix becomes zero and the dual pairings
are unchanged.  Thus $e,f$ span a hyperbolic space.  Do this on both
sides and match the two hyperbolic spaces by the nondegenerate
extension just proved.  This extends the original subspace isometry.
For $A=k\times k$, the Hermitian space is a perfect pairing of two
$k$-spaces; extending dual bases proves the same statement directly.
These arguments include the real unitary groups.

It follows locally that every tuple with Gram matrix $T$ is $g_vX_0$.
At almost all finite places the tuple and $X_0$ are integral and have
unimodular Gram matrix.  The isometry can then be chosen integral.
Here are the lifting details.  Over the residue field, GMR6's
orthogonal-complement argument supplies an isometry of the two frames.
The equation $g^*Jg=J$ is smooth over the valuation ring: its derivative
is onto the Hermitian matrices because two and $J$ are units.  Hence
the residue isometry lifts successively modulo $\varpi^m$ to an
integral unitary element.  If its frame is correct only modulo
$\varpi^m$, the next error $Z$ satisfies
$X_0^*JZ+Z^*JX_0=0$ modulo $\varpi$.  Relative to the nondegenerate
span of $X_0$ and its orthogonal complement, extend
$X_0\mapsto Z$ to a skew-adjoint endomorphism by choosing its other
off-diagonal block to be the negative adjoint and its final diagonal
block to be zero.  The correction $1+\varpi^m A$ is unitary to the
next order and corrects that frame error.  Smooth lifting corrects its
remaining unitary error at higher order.  Iteration converges in the
complete valuation ring.  Thus the local choices lie in
$U(J)(\mathcal O_v)$ almost everywhere and assemble to an actual
adelic unitary element.  This proves the orbit equality.

We also check its topology.  On the independent locus the tangent
space to the Gram fiber consists of the $Z$ with
$X^*JZ+Z^*JX=0$.  Any such partial map $X\mapsto Z$ extends to a
skew-adjoint endomorphism: extend the columns of $X$ to a basis, form
the sesquilinear matrix whose prescribed columns are
$h(-,Z)$, fill the prescribed rows by skew-adjointness, and put zero
in the unspecified diagonal block.  The condition on $Z$ is exactly
the compatibility on the overlapping prescribed block.  Multiplication
by $J^{-1}$ recovers the desired endomorphism.  The orbit differential
is therefore onto.  The inverse function theorem gives analytic local
sections and makes the local orbit map open.

At a good integral place the same argument works over the valuation
ring.  The compact fiber $C_v$ has a finite clopen cover on which there
are continuous sections with values in $U(J)(\mathcal O_v)$; a finite
disjoint refinement patches these sections to a continuous section
on $C_v$.  Products of these sections, with the finitely many other
local sections, give a continuous section on an adelic neighborhood.
Consequently the continuous bijection from the group quotient to
the fiber is an open map, hence a homeomorphism.

Finally every local $g_v$ preserves $\omega_{v,T}$ by GMR5, and at
almost every place it also preserves $C_v$.  Formula (GMR.16) on
cylinder sets proves invariance of $\nu_T$; those sets generate the
Borel sigma algebra.  No normalization scalar occurs in this action.
\end{proof}
% END GMR10

% BEGIN GMR11
\paragraph{GMR11. Quotient normalization for every section.}
\label{claim:GMR11}
Let a second countable locally compact group $G$ act transitively on
a locally compact space $Y$, let $y_0\in Y$ have closed stabilizer $H$,
and suppose $G/H\simeq Y$ with continuous local sections of the orbit
map.  Let $\rho$ be a nonzero invariant Radon measure on $Y$, positive
on nonempty open sets.  For every supplied left Haar measure $\mu_G$
there is a unique left Haar measure $\mu_H$ for which
\[
 \int_G f\,d\mu_G
   =\int_Y\int_H f(s(y)h)\,d\mu_H(h)\,d\rho(y)
 \tag{GMR.19}
\]
for all nonnegative measurable $f$ and all sections $s$, with no
measurability requirement on the section.  Measures supported on an
orbit in an ambient space satisfy the identical formula on that
ambient space by extension by zero.

\begin{proof}
Start with any left Haar measure $\eta$ on $H$.  For $f\in C_c(G)$,
the integral $q_f(y)=\int_H f(s(y)h)\,d\eta(h)$ is independent of $s$:
two representatives of the same coset $gH$ differ by right
multiplication of $s(y)$ by an element of $H$, and that change is a
\emph{left} change of variable in the inner integral.
Using a continuous local section shows that $q_f$ is continuous.
Indeed, on a compact neighborhood in its domain, the contributing
$h$ lie in the compact set formed from that section's image and
$\operatorname{supp}f$, intersected with $H$; dominated convergence
applies on this fixed compact set.  Its support is contained in the
compact image of $\operatorname{supp}f$ in $Y$.

Choose a countable cover by section domains and disjointify it to
obtain a Borel section.  For arbitrary nonnegative measurable $f$, its
inner integral is then measurable, and integrating it against $\rho$
defines a measure $L_\eta$ on $G$.  Countable additivity follows by
monotone convergence.  The preceding compact-support argument proves
local finiteness.  Positivity of both measures on neighborhoods gives
$L_\eta\ne0$.
For $g\in G$, write $g s(y)=s(gy)c(g,y)$ with $c(g,y)\in H$.
Left invariance of $\eta$ and invariance of $\rho$ give
$L_\eta(f\circ L_g)=L_\eta(f)$.  Thus $L_\eta$ is a left Haar measure
and $L_\eta=c\mu_G$ for some finite $c>0$.
Set $\mu_H=c^{-1}\eta$.  This proves (GMR.19).

Every other Haar measure on $H$ is a positive multiple of $\eta$;
a nonzero compactly supported test function in (GMR.19) determines
that multiple.  This proves uniqueness.  For any other section,
including a nonmeasurable one, the inner integral is pointwise the
same function of $y$ as for the Borel section.  Thus the original
all-section formula follows literally.  Values away from the orbit
are irrelevant when its complement has measure zero.
\end{proof}
% END GMR11

% BEGIN GMR12
\paragraph{GMR12. The full original global stabilizer-measure supplier.}
\label{claim:GMR12}
Let $E/F,c,n,J$ satisfy the original CM and nondegeneracy hypotheses.
Retain the original premises \texttt{SWa1} and
\texttt{StabilizerFirstChernCompatibility}.  For \emph{every} character
$\psi$ satisfying \texttt{IsCharAF}, every $a$ with $1\leq a$ and
$2a<n$, every rational ordered $a$-tuple $X_0$ whose Gram matrix
$T$ is nondegenerate, every witness that $T$ is Hermitian, and every
supplied Haar measure $\mu_G$ on $U(J)(\mathbb A_F)$, there exist
$\nu_T$ and a Haar measure $\mu_{X_0}$ on the actual pointwise
stabilizer of $X_0$ such that
\[
 \texttt{IsOmT}(\psi,T,\nu_T),\qquad
 \texttt{IsCompatT}(X_0,\mu_G,\mu_{X_0},\nu_T).
 \tag{GMR.20}
\]
This is the full statement of
\texttt{exists\_adelicTorusMeasure\_stabiliserHaar}.

\begin{proof}
Construct $\nu_T$ by GMR9.  GMR10 identifies its support with the
actual orbit, proves invariance, and provides the quotient topology
and local sections.  The group is a closed matrix group in the
adelic general linear group; it is locally compact and second
countable.  Its stabilizer is closed, since it is the zero set of
the continuous equations $gX_0=X_0$.  Apply GMR11 to the supplied
$\mu_G$, not to a newly chosen normalization.  Its conclusion is
exactly the all-section formula and support condition in
\texttt{IsCompatT}.  GMR9 supplies all the quantifiers in
\texttt{IsOmT}.  A witness of Hermitian symmetry merely specifies
the same element of the Hermitian subspace and changes none of these
measures.

This proves the original supplier under its exact premises.
The proof in fact did not use \texttt{SWa1} or
\texttt{StabilizerFirstChernCompatibility}: their role in the source
construction is replaced here by the explicit local construction and
the positive product in GMR6--GMR9.  This observation does not prove
either of those premises and does not assert any equality of
arithmetic quotient volumes.
\end{proof}
% END GMR12

\subsection{The local flag stabilizer and its product normalization}

Fix a standard rational frame
$P=(e_1,\ldots,e_r,f_1,\ldots,f_r,z)$, meaning
\[
 h(e_i,e_j)=h(f_i,f_j)=0,\quad h(e_i,f_j)=\delta_{ij},\quad
 h(e_i,z)=h(f_i,z)=0,\quad h(z,z)=-T_0,\quad T_0\in F^\times.
 \tag{GMR.21}
\]
Set $F_0=(e_1,\ldots,e_r,z)$ and
$T_F=\operatorname{diag}(0,\ldots,0,-T_0)$.
At a fixed finite place, let $S_0$ be the pointwise stabilizer of $F_0$
and let $S$ consist of the unitary elements acting upper
unitriangularly on this ordered tuple.  These are the literal
\texttt{S0} and \texttt{Ssub}, not conjugate replacement groups.
Let $\mathcal U=U_{r+1}(A_v)$ and identify its underlying space with
$A_v^{\{(i,j):i<j\}}$ by its strict upper entries.

% BEGIN GMR13
\paragraph{GMR13. The independent singular fiber is the local orbit.}
\label{claim:GMR13}
The tuple $F_0$ is independent and its Gram matrix is $T_F$.
The independent part of $\gamma^{-1}(T_F)$ is exactly $G_vF_0$.
Its canonical measure $\omega_{v,T_F}$ is a nonzero invariant Radon
measure on that orbit, and the orbit map has continuous local
sections.  For every supplied Haar measure $\mu_{G,v}$ there is a
Haar measure $\mu_0$ on the actual $S_0$ satisfying the full original
\texttt{IsCompatLoc} with $\omega_{v,T_F}$.
More generally, prescribe arbitrary additive Haar measures $dX,dT$
on the two spaces and any family $\mathrm{fib}$ satisfying their
original \texttt{IsFibreFamily} condition.  The same conclusion holds
with $\mathrm{fib}(T_F)$ in place of the canonical measure, for every
supplied $\mu_{G,v}$.  Thus the full original
\texttt{exists\_haar\_frameStabiliser\_compatible\_fibre} is included;
its two additive measures need not be self-dual.

\begin{proof}
If $\sum_i a_i e_i+bz=0$, pairing with the $f_j$ gives $a_j=0$ and
then pairing with $z$ gives $b=0$.  This proves independence, also
after scalar extension.  Its Gram matrix is (GMR.21).
For any other independent tuple of that Gram matrix, the last
vector is nonisotropic and its first $r$ vectors span a totally
isotropic subspace of its orthogonal complement.  The degenerate
subspace extension proved in GMR10 sends the ordered tuple $F_0$
to this tuple.  Thus the whole independent fiber is one actual orbit.

GMR4 makes its fiber measure Radon on the orbit and strictly positive
near every point; $F_0$ supplies a point, so it is nonzero.
GMR5 proves invariance.  The skew-adjoint extension of a tangent
map in GMR10 applies to any independent subspace, and therefore
gives local sections here as well.  The orbit is locally closed:
it is the closed Gram level set intersected with the open
independent locus.  Its stabilizer is closed and its quotient has
this orbit topology.  GMR11 now supplies $\mu_0$ relative to the
\emph{supplied} $\mu_{G,v}$ and proves the formula for every section.
GMR5 identifies the fiber measure with the source's
\texttt{omAt} at $T_F$.  This proves the original local
frame-stabilizer compatibility supplier even though $T_F$ is singular
when $r>0$.
For the more general assertion, perform the GMR4 chart construction
with exactly the given $dX,dT$.  The positive chart densities give
the same Radon and nonvanishing properties.  GMR5's pointwise
uniqueness identifies the given family with this constructed one,
and its determinant-one covariance gives unitary invariance for
these normalizations as well.  GMR11 now applies without any change.
\end{proof}
% END GMR13

% BEGIN GMR14
\paragraph{GMR14. Unitary Haar measure is unimodular.}
\label{claim:GMR14}
Every Haar measure on the actual local unitary group $G_v$ is both
left and right invariant.  In particular conjugation and right
translation by the elements of $S$ preserve the supplied Haar measure
on $G_v$.

\begin{proof}
The Lie algebra consists of $A_v$-linear endomorphisms $A$ with
$A^*J+JA=0$.  On it the trace pairing
$(A,B)\mapsto\operatorname{Tr}_{A_v/k_v}\operatorname{tr}(AB)$ is
nondegenerate.  To verify this, extend scalars from $k_v$ to the
quadratic etale algebra: the unitary Lie algebra becomes the full
matrix Lie algebra and this pairing becomes a nonzero scalar
multiple of its nondegenerate trace pairing.  Faithful scalar
extension then proves nondegeneracy over $k_v$ itself.  In the split
case this is also immediate from the description as
$\mathfrak{gl}_n(k_v)$.

Conjugation preserves this pairing, so
$(\det\operatorname{Ad}g)^2=1$.  Its absolute determinant is therefore
one.  In local analytic coordinates a Haar density may be formed
from a left-invariant top differential form.  Conjugation by $g$
multiplies that density by the absolute determinant of its
differential at the identity, namely
$|\det\operatorname{Ad}g|=1$.
Equivalently, comparing the translated coordinate densities proves
the familiar modular formula directly.  Hence conjugation preserves
Haar measure.  Composing conjugation with a left translation gives
a right translation, so right translations preserve it as well.
Haar uniqueness makes this conclusion independent of the initial
positive normalization.
\end{proof}
% END GMR14

% BEGIN GMR15
\paragraph{GMR15. Every element of the flag stabilizer preserves its kernel Haar measure.}
\label{claim:GMR15}
For every Haar measure $\eta$ on $S_0$ and every $s\in S$,
conjugation by $s$ preserves $\eta$.  In particular $S_0$ is
unimodular.  This proves the normalization assertion needed in the
original \texttt{similitude\_conj\_preserves\_unitary\_haar\_at\_finite\_place}
at the actual frame and for every Haar normalization.
In that original interface the additional prescribed additive Haar
measures and the prescribed fiber family may be any allowed ones;
the conclusion does not depend on them and our proof does not use
them.  Only the actual standard frame and the supplied Haar measure
on $S_0$ occur in the conclusion.

\begin{proof}
Let $u\in\mathcal U$ be the unique matrix with $sF_0=F_0u$.
The map $R_u:X\mapsto Xu$ commutes with the left $G_v$ action and
preserves the independent fiber of $T_F$, since $u^*T_Fu=T_F$.
It preserves $\omega_{v,T_F}$ by GMR5.  It also shows that $s$
normalizes $S_0$.

First use the measure $\mu_0$ fixed by GMR13's quotient formula.
If $C_s(h)=shs^{-1}$, its pushforward is another Haar measure,
so $(C_s)_*\mu_0=c\mu_0$ for a finite $c>0$.
Choose a section $q$ of the orbit map and let
$q'(y)=q(R_u^{-1}y)s$; this is again a section because
$sF_0=R_uF_0$ and $R_u$ commutes with the action.
Apply GMR13's formula to $f(\,\cdot\,s^{-1})$ using $q'$.
Its inner argument is
\[
 q'(y)h s^{-1}=q(R_u^{-1}y)\,s h s^{-1}.
\]
GMR14 says that the left integral is $\int f\,d\mu_{G,v}$.
Changing $y$ by $R_u$ and then using the Haar factor $c$ says that
the right integral is $c\int f\,d\mu_{G,v}$.
Choose a nonnegative compactly supported $f$ with finite positive
integral.  We conclude $c=1$.

Every other Haar measure $\eta$ is a positive scalar multiple of
$\mu_0$, so it is preserved as well.  Taking $s\in S_0$ proves
invariance under inner automorphisms of $S_0$.  Together with left
invariance this is right invariance, proving unimodularity.
Thus the product construction below uses an established exact
normalizer action, rather than an unspecified modular factor.
\end{proof}
% END GMR15

% BEGIN GMR16
\paragraph{GMR16. The actual upper-unitriangular quotient and a continuous section.}
\label{claim:GMR16}
The map $\pi:S\to\mathcal U$ determined by $sF_0=F_0\pi(s)$ is a
continuous surjective homomorphism with kernel $S_0$.  It has a
continuous section $t$.  Consequently
\[
 S_0\times\mathcal U\longrightarrow S,\qquad (s_0,u)\longmapsto s_0t(u)
 \tag{GMR.22}
\]
is a homeomorphism.  The source measure
$\texttt{sdE}(\psi,v,\texttt{UIdx}(r))$ is a left and right Haar measure
on $\mathcal U$ under its strict-upper-coordinate identification.

\begin{proof}
Independence of $F_0$ makes its coefficients unique; taking products
proves the homomorphism assertion and taking coordinates proves
continuity.  Its kernel is the pointwise stabilizer, exactly $S_0$.
To prove surjectivity with the needed topology, write
\[
 u=\begin{pmatrix}A&b\\0&1\end{pmatrix},
 \qquad A\in U_r(A_v),\quad b\in A_v^r.
\]
Use the frame's nondegenerate subspace spanned by $e,f,z$, and let the
following map act as the identity on its orthogonal complement:
\[
\begin{split}
 t(u)e&=eA,\\
 t(u)z&=z+eb,\\
 t(u)f&=\left(f+T_0^{-1}z b^*
                  +(2T_0)^{-1}e b b^*\right)A^{-*}.
\end{split}
 \tag{GMR.23}
\]
Here $e,f$ denote the matrices with those columns and
$A^{-*}=(A^*)^{-1}$.  Formula (GMR.21) gives
$h(t e,t f)=1$, $h(tz,t f)=0$ and $h(t f,t f)=0$:
in the last identity the two terms with coefficient $(2T_0)^{-1}$
cancel the term $-T_0\cdot T_0^{-2}$ from the $z$ coordinate.
The other pairings are unchanged.  Thus $t(u)$ is unitary and has
the prescribed action on $F_0$.  Its formulas are continuous in
$A,b$, including $r=0$, when the section is the identity.

Every $s\in S$ has the unique expression
$(s t(\pi(s))^{-1})t(\pi(s))$.  Its first factor belongs to $S_0$,
and these inverse formulas are continuous, proving (GMR.22).

Finally in the strict-upper coordinates the product law is
\[
 (xy)_{ij}=x_{ij}+y_{ij}+\sum_{i<k<j}x_{ik}y_{kj}.
 \tag{GMR.24}
\]
For fixed $x$, this is an affine triangular change of the $y$
coordinates with diagonal entries one, ordered by increasing
$j-i$; the same holds with the roles reversed.  It preserves
additive Haar measure by repeated one-coordinate translations and
Fubini, or by its determinant one.  GMR2 identifies the source
\texttt{sdE} with a nonzero additive Haar measure on this coordinate
space.  Therefore it is also left and right Haar for the group law
(GMR.24), with exactly that self-dual normalization.
\end{proof}
% END GMR16

% BEGIN GMR17
\paragraph{GMR17. The full all-section product formula on the flag stabilizer.}
\label{claim:GMR17}
Let $\mu_0$ be the measure on $S_0$ from GMR13, and let $\mu_U$ be
the literal \texttt{sdE} measure in GMR16.  There is a Haar measure
$\mu_S$ on the actual $S$ for which
\[
 \int_S f(s)\,d\mu_S(s)
 =\int_{\mathcal U}\int_{S_0}f(s_0t'(u))\,d\mu_0(s_0)\,d\mu_U(u)
 \tag{GMR.25}
\]
for every nonnegative measurable function on the ambient unitary
group and every map $t'$ inducing the prescribed upper-unitriangular
action on $F_0$.  In particular this is the exact original
\texttt{IsProdU}, with no measurability assumption on $t'$.

\begin{proof}
Push forward $\mu_0\otimes\mu_U$ by (GMR.22).  This defines a
nonzero Radon measure on $S$, finite on compact sets and positive on
open sets.  To prove left invariance, fix $g\in S$ and put
$v=\pi(g)$.  There is a factor $c(g,u)\in S_0$ such that
$g t(u)=c(g,u)t(vu)$.  Thus
\[
 g s_0 t(u)=(g s_0 g^{-1})c(g,u)t(vu).
\]
GMR15 permits the first conjugation in the inner integral, and the
unimodularity of $S_0$ permits right multiplication by $c(g,u)$.
GMR16 permits the change $u\mapsto vu$ in the outer integral.
The resulting integral is the original one.  Hence the pushed-forward
measure is left invariant and is Haar.

If $t'$ has the same frame action as $t$, then
$t'(u)t(u)^{-1}\in S_0$.  Right invariance of $\mu_0$ makes the inner
integrals with $t'$ and $t$ equal pointwise in $u$.  The one with the
continuous section $t$ is measurable by Tonelli.  Thus the formula
also holds for an arbitrary $t'$; its inner integral is the same
measurable function.  Pushing $\mu_S$ to the ambient group gives
the source's \texttt{amb} convention in \texttt{IsProdU} exactly.
\end{proof}
% END GMR17

% BEGIN GMR18
\paragraph{GMR18. The full original normalized local frame-stabilizer supplier.}
\label{claim:GMR18}
Assume the original CM, Hermitian, \texttt{SWa1} and
\texttt{StabilizerFirstChernCompatibility} premises.  Let $\psi$
satisfy \texttt{IsCharAF} and \texttt{IsStdChar}.  For every finite
place $v$, every $r$, every actual standard frame $P$ with
$2(r+1)<n$, and every supplied Haar measure $\mu_{G,v}$, there is a
Haar measure $\mu_{S,v}$ on the actual $\texttt{Ssub}(P)$ such that
\[
 \texttt{LocTieW}(c,J,\psi,v,P,\mu_{G,v},\mu_{S,v}).
 \tag{GMR.26}
\]
This is the full original
\texttt{exists\_frameStabilizerHaar\_normalised\_fibreMeasure}.

\begin{proof}
Use GMR2--GMR5 at $a=r+1$ to obtain the source's selected self-dual
measures and its selected entire fiber family.  Since $T_F$ is
Hermitian, \texttt{omAt} is that family's value at $T_F$, not its
zero fallback.  GMR13 gives a Haar measure $\mu_0$ on the literal
pointwise stabilizer $S_0$ with \texttt{IsCompatLoc} relative to the
supplied $\mu_{G,v}$.  GMR2 and GMR16 show that the selected
\texttt{sdE} exists and is the required Haar measure in the
upper-unitriangular coordinates.  GMR15 and GMR17 construct $\mu_{S,v}$
and prove the all-section \texttt{IsProdU} formula.

These are exactly the three entries of \texttt{LocTieW}: the witness
$\mu_0$ is Haar, is compatible with \texttt{omAt}, and gives the
product formula with \texttt{sdE}.  All the ambient groups, the
character, the frame and the supplied $\mu_{G,v}$ have been kept.
The unused original volume premises are retained in the statement;
no local density or normalization supplier has been assumed in their
place.  The argument even allows $r=0$: $\mathcal U$ is the trivial
group with self-dual mass one, and (GMR.25) then gives
$\mu_S=\mu_0$.
\end{proof}
% END GMR18

\subsection{Simultaneous rescaling without changing the supplied bridge}

% BEGIN GMR19
\paragraph{GMR19. The full original \texttt{RescaleStep3}.}
\label{claim:GMR19}
Assume the original \texttt{SiegelWeilUnitary}\allowbreak
\texttt{Convergent} premise.
For every original $D$, every supplied $X:\texttt{HidaSetting}(D)$,
every supplied $h_b:\texttt{AdBridgeIso}(X.B.I)$, and every witness
of $X.\texttt{MeasTie}(h_b)$, there are normalized measure records
$m_G$ and $m_W(u)$ such that
\[
 X'=\texttt{withMeas}(X,m_G,m_W)
 \tag{GMR.27}
\]
has all the following properties simultaneously:
\[
\begin{gathered}
 X'.\texttt{MeasTie}(h_b),\qquad
 \exists\psi:\mathbb A_F\longrightarrow\mathbb C^\times,\quad
 \texttt{IsCharAF}(\psi)\ \wedge\ \texttt{IsStdChar}(\psi),\\
 \texttt{MeasCompat\ensuremath{\psi}}(X',\psi),\qquad
 \forall w\ \texttt{LocCompatX}(X',\psi,w).
\end{gathered}
 \tag{GMR.28}
\]
Here $F=X.B.I.F$ and every group, frame and tensor coordinate is the
one belonging to $X$.  We can, and do, leave
$\mu_G,\mu_{G,\infty},\mu_{G,f},\mu_{G,v}$ unchanged and retain
\[
 m_W(u).\mu_{N,\infty}=X.\texttt{measW}(u).\mu_{N,\infty},\qquad
 m_W(u).\mu_{H,\infty}=X.\texttt{measW}(u).\mu_{H,\infty}
 \tag{GMR.29}
\]
for every word.  The same single $\psi$ works for all words, all
finite places, and all auxiliary tuples of length $2D.R.rS+2$.

\begin{proof}
Put $r=D.R.rS$, $n=X.B.I.n$.  The setting's CM and Hermitian fields
give the hypotheses used above.  Its rank equation is
$n=D.R.k+D.R.m-1$.  Write $D.R.m=2(r+1)$, using evenness and
$D.R.m\geq2$.  Since $D.R.k\geq D.R.m+3$, we have
$n\geq4r+6$, in particular the source bound $4r+5\leq n$.
Thus
\[
 2(2r+2)<n,\qquad 2(2r+1)<n,\qquad 2(r+1)<n.
 \tag{GMR.30}
\]

For every word $u$, let $P_u$ be the source's actual
$\texttt{frameOf}(X,u)$.  The setting's basis and its equality for
$J$ give
\[
 h_J(\operatorname{crd}x,\operatorname{crd}y)
   =X.\iota_E^{-1}(h_{VS}(x,y)).
 \tag{GMR.31}
\]
Indeed expand both vectors in that basis, use its $J$ entries, and
use that $X.\iota_E$ respects the involutions.  The stored
$ee,ff,ef,zu\_orth,zu\_norm$ equations therefore give (GMR.21)
for $P_u$, with $T_u$ the nonzero rational integer specified by the
word.  It is nonzero because that integer is prime in the original
word datum.  The ordered tuple
$\texttt{tupH}(P_u)=(e_1,\ldots,e_r,z,f_1,\ldots,f_r)$ has
nondegenerate Gram matrix: after moving $z$ last it is
\[
 \begin{pmatrix}0&1&0\\1&0&0\\0&0&-T_u\end{pmatrix},
 \tag{GMR.32}
\]
where the first two diagonal blocks have size $r$.  Its determinant
is $(-1)^r(-T_u)\ne0$.  This checks the original standard-frame
and rank suppliers directly on $X$.

Choose \emph{now}, once and for all, the standard character of GMR1
on $\mathbb A_F$.  No word, place, or auxiliary tuple has been chosen
at this stage.  Specializing the retained Siegel--Weil premise at
the actual $F,E,c,n,J$ gives its finite-volume and smooth-density
components.  By their literal definitions these are the source's
\texttt{SWa1} and \texttt{StabilizerFirstChernCompatibility} premises
in GMR12 and GMR18.  We use those original statements at this scope;
their constructive proofs above did not otherwise require these
volume components.

First handle all auxiliary tuples $X_0$ of length $2r+2$.
If $X_0$ has nondegenerate Gram matrix, apply GMR12 with the
unchanged global measure $X.\texttt{measG}.\mu_G$.
Choose its stabilizer Haar measure $\mu'_B(X_0)$ and fiber measure
$\nu_{X_0}$.  If its Gram matrix is degenerate, retain the existing
$X.\texttt{measG}.\mu_B(X_0)$; the source imposes no
\texttt{MeasCompat\ensuremath{\psi}} fiber condition on that tuple.
Thus a single choice function $\mu'_B$ is Haar for every tuple and
has the required compatibility for every nondegenerate one.

Next fix a word $u$ temporarily.  Apply GMR12 to the actual tuple
$\texttt{tupH}(P_u)$, of length $2r+1$, using (GMR.30)--(GMR.32).
This gives a global Haar measure $\mu_{H,u}$ and its global fiber
measure, compatible with the same $\mu_G$ and the same $\psi$.
For every finite place $v$, apply GMR18 to $P_u$ and the unchanged
supplied $X.\texttt{measG}.\mu_{G,v}$; choose the resulting
$\mu_{S,u,v}$ with its full \texttt{LocTieW}.
Apply AMD17 with precisely the old $\mu_{N,\infty}$ and
$\mu_{H,\infty}$ for this word, this newly fixed $\mu_{H,u}$,
and this family $\mu_{S,u,v}$.  It produces a normalized
$m_W(u)$ retaining all these entries and with
$\texttt{covol}(m_W(u).\mu_N)=1$.

To make the normalization transparent, AMD16 first chooses the
global $N$-measure of literal covolume one.  AMD13 absorbs the
two required global-to-infinite Haar scalars into $\mu_{N,f}$ and
$\mu_{H,f}$ separately.  AMD7 supplies the finite $S$-measure as
the product pushforward of those two finite factors.  AMD14 verifies
all eleven \texttt{MeasW.IsTextNorm} clauses.  In particular neither
measure in (GMR.29), the prescribed global $H$-measure, nor any
prescribed local $S$-measure is rescaled.  The source record does
not require an additional equality identifying this finite
$S$-measure with a restricted product of the local $S$-measures;
no such extra equality is used here.
Choosing these witnesses for all words gives the required family
$m_W$.  All choices have occurred after the common $\psi$ was fixed.

Define $m_G$ by changing only $\mu_B$ to $\mu'_B$ in the old
\texttt{MeasG} record.  Its Haar clauses hold and its product
splitting clause is unchanged, so it satisfies the full
\texttt{MeasG.IsTextNorm}.  Form $X'$ by (GMR.27).
The first three entries of \texttt{MeasTie} are unchanged: they
refer to the same $\mu_G,\mu_{G,\infty},\mu_{G,f}$, the same
block measure $X.B.\mu_f$, and the same maps through $h_b$.
Its last two entries, for each word, are the ambient pushforwards
of $\mu_{N,\infty}$ and $\mu_{H,\infty}$; (GMR.29) makes those
entries unchanged as well.  Thus the exact supplied bridge $h_b$
still witnesses $X'.\texttt{MeasTie}$.

Finally unfold \texttt{MeasCompat\ensuremath{\psi}}.  Its first clause is,
for each word, the global $H$ fiber compatibility together with
unipotent covolume one; these were just supplied.  Its second clause
is global compatibility for every nondegenerate auxiliary tuple;
that is the choice of $\mu'_B$.  Unfolding
\texttt{LocCompatX} gives \texttt{LocTieW} at every finite place
for each word, exactly the local choices above.  We have therefore
proved every entry of (GMR.28), with the original order and range
of quantifiers.  This is the full original \texttt{RescaleStep3}.
\end{proof}
% END GMR19

\subsection{Exact next use and retained global analytic inputs}

GMR19 applies directly to the supplied-form core setting of BQS23,
which already consumes ASG, AMD and the completed Bessel chain.
It adds \texttt{MeasCompat\ensuremath{\psi}} and all \texttt{LocCompatX} to
that \emph{same} setting while preserving its bridge and its
infinite measures.  It does not construct the original supplied
\texttt{Forms0} datum.  The source's stronger all-family
\texttt{exists\_setting} subsequently uses the entire family
$F(u)$ to produce \texttt{Transport} and \texttt{LevelMatch};
that step has distinct genus, level-conjugacy and representation
intertwiner producers.  If such a transport is already supplied,
\texttt{trMeas} keeps its Witt map, adelic algebra map, representatives
and representation intertwiner literally unchanged under
\texttt{withMeas}.  No new transport existence is inferred from a
measure normalization.

For clarity, the retained source premise
\texttt{SiegelWeilUnitaryConvergent} says the following for every CM
extension and every nondegenerate Hermitian ambient space.  Haar
measures and arithmetic fundamental domains exist, and each such
fundamental domain has finite Haar mass.  For every tuple length
$a$ with $2a<n$, every nondegenerate rational Hermitian $T$, every
rational coefficient basis, every factorizable test function whose
infinite part is smooth, compactly supported and rapidly decreasing,
and every pure finite Schwartz--Bruhat factorization, all the source
Haar, fundamental-domain, character and self-dual-measure choices
admit continuous Gram-pushforward densities on the nondegenerate
Hermitian locus.  Their values at $T$ are multipliable, and their
infinite value times the finite product equals the normalized
arithmetic theta mean.  The source defines that mean as the complex
integral divided by the real mass of the chosen fundamental domain.
This exact universal premise is retained; the local disintegration
and positive product proved here are not an assertion of that
arithmetic theta identity.

In particular neither equality of the arithmetic stabilizer and
ambient volumes, nor the later orbital-counting identity, follows
just by replacing its Haar measures.  The next original consumer
\texttt{inVolumes\_of\_compat} still takes its named
\texttt{InVolumesAt} input.  The actual
\texttt{CountIdempotentsClass}, local Bessel uniqueness, the
all-family column transports, and the final global spectral
identifications remain separate source obligations.  This chapter
closes the measure-rescaling supplier and its two analytic measure
constructors, with their complete original scopes, rather than a
whole selected global obligation.
